\chapter{Разработка архитектуры системы управления динамической платформой} \label{chapt5}

\section{Структурная схема ПАК управления динамической платформой} \label{sect5_1}

Описанные в главе \ref{chapt4} алгоритмы реализуются в составе программно-аппаратного комплекса управления СП, который содержит следующие основные компоненты, показанные на рисунке \ref{img:MPStruct}:

\begin{itemize}
	\item Контроллер верхнего уровня / контроллер движения СП реализует алгоритмы планирования траектории и комплексного управления приводами СП,
	\item Контроллер привода осуществляет управление локальной степенью подвижности платформы
	\item Клиент системы подвижности -- устройство, подключенное к сети Ethernet, являющееся задатчиком команд для сервера системы подвижности,
	\item Мотор-редукторы -- исполнительные устройства СП, осуществляют позиционирование выходного звена редуктора (кривошипа),
	\item Датчики угла -- основной источник данных о положении платформы, необходимый для нормального функционирования приводов СП,
	\item MEMS-датчик движения (совмещенный 3DOF акселерометр и гироскоп), используется для проверки геометрии СП, контроля состояния датчиков угла.
\end{itemize}

\begin{figure}[!htbp]
    \center{
        \hfill
        \subcaptionbox[List-of-Figures entry]{\label{img:MPStructCommonCentralized}} {\includegraphics[width=0.5\linewidth]{MPStructCommonCentralized}}%
        \hfill       
        \subcaptionbox{\label{img:MPStructCommon}} {\includegraphics[width=0.5\linewidth]{MPStructCommon}}
        \hfill
    }

    Варианты компоновки ПАК СП
    \caption{Структурная схема ПАК СП} % Этот текст попадает в названия рисунков в списке рисунков
    \label{img:MPStruct}
\end{figure}

Как показано на рисунке \ref{img:MPStruct} существует два подхода к построению ПАК управления платформой: централизованная система управления \ref{img:MPStructCommon} и децентрализованная структура \ref{img:MPStructCommon}. Первая основана на использовании последовательных/аналоговых интерфейсов для управления мотор-редукторами и считывания данных с датчиков угла. Такой вариант предполагает сведение всех линий к единому преобразователю интерфейсов, что увеличивает количество используемых кабелей. Второй вариант предполагает использование промышленной шины (CAN, RS485, Profibus, пр.) для объединения устройств. На шине задается адресация  и назначается ведущее (master) устройство, которым в случае ПАК СП является КД. При использовании шинной архитектуры упрощается подключение новых устройств. Так, для добавления нового контроллера привода в сеть достаточно организовать ответвление, подключить новое устройство и задать его уникальный адрес.

В практической реализации ПАК СП был выбран вариант \ref{img:MPStructCommon}, как более эффективное и масштабируемое решение построения системы управления СП. 

\section{Датчики активной безопасности системы подвижности. Алгоритмы обработки сигналов датчиков безопасности} \label{sect5_2}

Для обеспечения безопасной эксплуатации СП в составе ПАК используются т.н. датчики активной безопасности, схема подключения которых изображена на рисунке \ref{img:AdditionalSensors}.

\begin{figure}[!htbp]
	\center
	\includegraphics[width=\linewidth]{AdditionalSensors}
	\caption{Схема подключения дополнительных датчиков безопасности СП}
	\label{img:AdditionalSensors}
\end{figure}

Конвертер сигналов КС-01 предназначен для преобразования сигналов с датчиков типа <<сухой контакт>>, установленных на рабочем месте обучаемого в пакеты данных, отправляемых по CAN-шине по заданному алгоритму (по запросу, при изменении данных, с заданной периодичностью). Конвертер оснащен светодиодной индикацией, отображающей статус подачи питания, и передачу данных на CAN-шину. При размыкании контакта на одном из датчиков производится аварийный останов системы подвижности.

Двухрежимная система измерения отклонения (инклинометр) применяется для калибровки системы и используется в составе дополнительного контура безопасности СП. Прибор оснащен разъемом Molex Brad Microchange M12 для подключения к внутренней CAN-шине ПАК СП. 

В процессе работы при размыкании контактов одного из датчиков безопасности происходит немедленная аварийная остановка СП. Для корректного возобновления работы используется алгоритм, работа которого показана на рисунке \ref{img:DryContactAlgo}:

\begin{figure}[!htbp]
\center
\caption{К описанию алгоритма остановки и возобновления движения платформы}
\label{img:DryContactAlgo}
\includegraphics[width=\linewidth]{DryContactAlgo}
\end{figure}

\section{Архитектура встраиваемого ПО контроллера движения АТМ12} \label{sect5_3}
В состав ПАК СП входит встраиваемое ПО АТМ12 собственной разработки. ПО состоит из основной программы с интерфейсной частью и динамически подгружаемых модулей. На этапе загрузки основной программы (ОП) производится чтение параметров из конфигурационных файлов и подгружаются соответствующие модули. При этом в ОП есть возможность выбрать необходимые для загрузки модули и настроить соответствующие им файлы конфигурации <<на лету>>, без перезапуска ОП.

В ПО введены следующие разновидности подключаемых модулей:
\begin{itemize}
	\item управление подвижностью (интерфейс IPlatformFace),
	\item внешнее управление (интерфейс IManipulatorFace),
	\item математический модуль контроля за передаваемыми значениями (интерфейс IRestrictFace),
	\item дополнительные интерфейсы, подгружаемые в зависимости от обстоятельств (ICanFace);
\end{itemize}

Каждый из перечисленных типов модулей имеет свой виртуальный интерфейс, общий для всех однотипных библиотек с набором методов, которыми оперирует ОП при работе ПО. Библиотеки независимы друг от друг и, также как и ОП, могут быть использованы самостоятельно.

Диаграмма классов ОП представлена на рисунке \ref{img:UMLschemeCore}.

\begin{figure}[!htbp]
	\center
	\includegraphics[width=\textwidth]{UMLschemeCore}
	\caption{Диаграмма классов основной программы}
	\label{img:UMLschemeCore}
\end{figure}

Класс CNDPlatform представляет собой <<ядро>> ОП, агрегирующее объекты классов PlatformController и ManipulatorController, которые используются для управления платформой. CNDPlatform также реализует методы обработки событий сигналов графического интерфейса пользователя и хранение используемых конфигураций СП в виде списков List<ItemConfigInfo>, где ItemConfigInfo -- структура данных, содержащая такие параметры, как пути к библиотекам управления и ограничений, пути к конфигурационным файлам и т.д.

Класс PlatformController инкапсулирует взаимодействие с подгружаемыми модулями управления платформой и ограничений, осуществляет загрузку соответствующей разделяемой библиотеки (методы loadPlatformLibrary, loadRestrictLibrary) и проверяет ее корректность (методы isPlatformLibOk, isRestrictLibOk).

Интерфейс IPlatformFace включает реализуемые в подключаемых библиотеках функции управления положением платформы (setPosition, moveToDefault, setLegLength, stopAction), чтения состояния (getPosition, getLegLength, getPlatformState) и калибровки исполнительных систем(calibrate).

Интерфейс IRestrictFace включает функции проверки задаваемого положения платформы (checkPosition) и получения пределов по координатам (getMinPosition, getMaxPosition), которые рассчитываются индивидуально для каждого типа платформы.

Класс ManipulatorController инкапсулирует передачу данных от внешнего задатчика команд через ядро программы классу PlatformController. Интерфейс IManipulatorFace содержит функции setOperatingFunction и operateFunc, параметрами которых являются номер и указатель на функцию, определенную в IPlatformFace, что позволяет использовать в подключаемом модуле внешнего управления необходимые функции управления платформой.

На рисунке \ref{img:UMLschemeManipulator} изображено взаимодействие классов модулей внешнего управления с ОП. 

\begin{figure}[!htbp]
	\center
	\includegraphics[width=\textwidth]{UMLschemeManipulator}
	\caption{Диаграмма классов для модулей внешнего управления}
	\label{img:UMLschemeManipulator}
\end{figure}

Класс EthernetManip реализует функционал TCP/IP-сервера, принимающего подключения по указанному  порту mServPort. Согласно протоколу обмена данными с серверной частью ПО управления СП, EthernetManip ожидает команды на инициализацию системы подвижности, активацию приводов и перевод в рабочее положение. После подтверждения начала работы сервер в непрерывном режиме отрабатывает команды по положениям, скоростям и ускорениям согласно алгоритму, описанному в приложении \ref{list:restrict}.
Для работы протокола реализованы следующие обработчики событий: подключение клиента (onUserConnected), получения очередного пакета (onDataReceived), штатное/нештатное отключение клиента (onDisconnected), обработки ошибки (onError).

Класс CANManip функционирует схожим с EthernetManip образом, однако прием команд от задатчика производится по общей или выделенной сети CAN. Для передачи данных используется вспомогательный класс CDevCan. CDevCan также используется для обмена данными со вспомогательными датчиками (ИНК-01, КС-01). 

На рисунке \ref{img:UMLschemePlatforms} показано взаимодействие классов модулей управления платформой подвижности с ОП.

Были реализованы библиотеки управления следующими типами СП:
\begin{itemize}
	\item трехстепенной электромеханической подвижностью (Electro3D),
	\item двухстепенной электромеханической подвижностью (Electro2D),
	\item трехстепенной подвижностью на сервоприводах (Servo3D),
	\item шестистепенной пневматической подвижностью (Pneumo6D),
	\item шестистепенной подвижностью на актуаторах (Actuator6D);
\end{itemize}

Каждый из вышеперечисленных классов реализует виртуальные методы setPosition, getPosition, moveToDefault, calibrate и другие базовые функции управления платформой. При загрузке библиотеки управления платформой производится инициализация аппаратной части, запускается поток (или несколько потоков, в зависимости от задачи), в котором циклически производится опрос датчиков, расчет управляющих воздействий, обмен данными с исполнительными устройствами. Функция calibrate() используется для проведения автоматической калибровки исполнительного устройства с сохранением необходимой информации в лог-файлах.

Для работы с протоколами высокого уровня поверх CAN (CANOpen) в ПО выделен отдельный интерфейс ICanFace, реализующий такие функции, как setParameter(), updateParamSet(), setPDOMode(), proceedHoming() для работы с протоколами PDO и SDO. Функция setPosition() реализует процедуры

\begin{figure}[!htbp]
	\center
	\includegraphics[width=\textwidth]{UMLschemePlatforms}
	\caption{Диаграмма классов для модулей управления платформой}
	\label{img:UMLschemePlatforms}
\end{figure}
